\documentclass{rapportENSIAS}
\usepackage{lipsum}
\usepackage{minted}
\definecolor{bg}{rgb}{0.95,0.95,0.95} % Adjust the RGB values as needed
\title{Rapport ENSIAS - Template} %Titre du fichier

\begin{document}

%----------- Informations du rapport ---------

\titre{Titre du rapport} %Titre du fichier .pdf
%\UE{UE PRO} %Nom de la UE
\sujet{Compte Rendu des Travaux Pratiques} %Nom du sujet

\enseignant{Prénom \textsc{Nom}} %Nom de l'enseignant

\eleves{Rabii \textsc{ABOUAISSA} \\
		Prénom \textsc{Nom} } %Nom des élèves

%----------- Initialisation -------------------
        
\fairemarges %Afficher les marges
\fairepagedegarde %Créer la page de garde
\tabledematieres %Créer la table de matières


%------------ Corps du rapport ----------------
\section{Première section} 

\lipsum[3-4]%Effacer cette ligne et écrire le texte souhaité

\subsection{Subsection}

\lipsum[3-4] %Effacer cette ligne et écrire le texte souhaité

\section{Deuxième section}

\lipsum[3-5] %Effacer cette ligne et écrire le texte souhaité




%------------- Commandes utiles ----------------
\newpage %pour ecrire dans une nouvelle page
\section{Quelques commandes}

Voici quelques commandes utiles :

%------ Pour insérer et citer une image centralisée -----
\subsection{Commande pour insérer et citer une image centralisée}

\insererfigure{logos/ensias.png}{3cm}{Légende de la figure}{Label de la figure}
% Le premier argument est le chemin pour la photo
% Le deuxième est la hauteur de la photo
% Le troisième la légende
% Le quatrième le label
Ici, je cite l'image \ref{fig: Label de la figure}


%------- Pour insérer et citer une équation --------------
\subsection{Commande pour insérer et citer une équation}

\begin{equation} \label{eq: exemple}
\rho + \Delta = 42
\end{equation}

L'équation \ref{eq: exemple} est cité ici. 

% ------- Pour écrire des variables ----------------------

Pour écrire des variables dans le texte, il suffit de mettre le symbole \$ entre le texte souhaité comme : constante $\rho$. 

%------- Pour ecrire un code ------------------------------
\subsection{Commande pour ecrire un code}

%linenos pour avoir les nombres des lignes 
%specifier la language de programmation utilisé

\begin{listing}[!ht]
\begin{minted}[bgcolor=bg,linenos]{C}
#ifndef DECLARATION_H_INCLUDED
#define DECLARATION_H_INCLUDED

//declaration de structure qui englobe les donnees d'un adherant
typedef struct _donnee
{
    char code[20];
    char nom[20];
    char domaineInteret[20];
}donnees;

#endif // DECLARATION_H_INCLUDED
\end{minted}
\caption{exemple de caption de code C}
\label{listing 1}
\end{listing}

One-line code formatting also works with \texttt{minted}. For example, a small fragment of HTML like this:
\mint{html}|<h2>Something <b>here</b></h2>|
\noindent can be formatted correctly.


%-------- si vous voulez ajouter liste des codes sources de document
%----- un-comment these two lines
%\renewcommand\listoflistingscaption{Liste des codes source}
%\listoflistings

\end{document}
